% SGA 2，原 PDF 1--12。原目录 PDF 5--7 由分册/总本主文件重建。
% 原目录的逐项中文题名与原版页码见 work/sga2-00-frontmatter/contents-translation.json。
\SGAfrontchapter{原书题名与出版说明}
\SourcePage{1}
\begin{center}
布尔玛丽\Term{algebra}{代数}几何讨论班\\
1962 年

\bigskip
{\Large\bfseries \Term{coherent-sheaf}{凝聚层}的\Term{local-cohomology}{局部上同调}\\及\Term{local-scheme}{局部}与整体的 Lefschetz 定理}

\medskip
（SGA 2）

\bigskip
Alexander Grothendieck\\
（由一组听众整理）

\medskip
增补 Michèle Raynaud 女士的一讲

\bigskip
本书为《纯数学高等研究》（\textit{Advanced Studies in Pure Mathematics}）\\
第 2 卷的重排注释版；原卷由 North-Holland Publishing Company 于 1968 年出版。

\medskip
\texttt{arXiv:math/0511279v1 [math.AG]}，2005 年 11 月 10 日。
\end{center}
\SourcePage{2}
% 原 PDF 第 2 页为空白，译本不另造空白页。

\SGAfrontchapter{重排版前言（法文版）}
\SourcePage{3}
本书是 A. Grothendieck 等人所著《\Term{coherent-sheaf}{凝聚层}的\Term{local-cohomology}{局部上同调}及\Term{local-scheme}{局部}与整体的 Lefschetz 定理（SGA 2）》的重排注释版。原书为 \textit{Advanced Studies in Pure Mathematics} 第 2 卷，由 North-Holland Publishing Company 于 1968 年在阿姆斯特丹出版。这是 B. Edixhoven 发起的重版工程的第二部分；他已经重新出版了 SGA 1。本版重现原文，同时作了两类补充：一类是纠正排印错误的少量形式修改，另一类是称为“编者注”（N.D.E.）的页下注，用以说明原版中所提出问题的现有研究状况。此外，我们还在若干地方为证明增添了一些细节。为避免混淆，原注仍保留星号编号，新注则用整数编号。原版页码标在页边。

感谢承担了绝大部分初始 \LaTeX2e{} 录入工作的数学家：L. Bayle、N. Borne、O. Brinon、J. Buresi、M. Chardin、F. Ducrot、P. Graftiaux、F. Han、P. Karwasz、L. Koelblen、D. Madore、S. Morel、D. Naie、B. Osserman、J. Riou 和 V. Sécherre；也感谢 C. Sabbah 将本文调整为法国数学会（SMF）的版式。此外，感谢 J.-B. Bost、P. Colmez、O. Gabber、W. Fulton、S. Kleiman、F. Orgogozo、M. Raynaud 和 J.-P. Serre 所提出的意见和建议。

\begin{flushright}编者：Yves Laszlo\end{flushright}

\SGAfrontchapter{重排版前言（英文版）}
\SourcePage{4}
本书是 A. Grothendieck 等人所著《\Term{coherent-sheaf}{凝聚层}的\Term{local-cohomology}{局部上同调}及\Term{local-scheme}{局部}与整体的 Lefschetz 定理（SGA 2）》的更新重版。原书为 \textit{Advanced Studies in Pure Mathematics} 第 2 卷，由 North-Holland Publishing Company 于 1968 年在阿姆斯特丹出版。这是 B. Edixhoven 发起的 SGA 工程的第二部分；他曾为 SGA 1 编印新版。本版旨在重现原文，并作若干修改，其中包括少量排印错误的订正，以及说明初版所提问题当前状况的编者页下注（N.D.E.）。我们也为某些证明补充了更多细节。为避免可能的混淆，原脚注以星号编号，新脚注则以整数编号。原版页码写在正文页边。

感谢承担了绝大部分初始 \LaTeX2e{} 排版工作的数学家：L. Bayle、N. Borne、O. Brinon、J. Buresi、M. Chardin、F. Ducrot、P. Graftiaux、F. Han、P. Karwasz、L. Koelblen、D. Madore、S. Morel、D. Naie、B. Osserman、J. Riou 和 V. Sécherre；也感谢 C. Sabbah 将本文调整为 SMF 的版式。还要感谢 J.-B. Bost、P. Colmez、O. Gabber、W. Fulton、S. Kleiman、F. Orgogozo、M. Raynaud 和 J.-P. Serre 所提出的意见和建议。

\begin{flushright}编者：Yves Laszlo\end{flushright}

% 原目录由共享主文件根据中文标题和译本页码重建，不在正文模块插入全书目录。
\SourcePage{5}
\SourcePage{6}
\SourcePage{7}
\SourcePage{8}
% 原 PDF 第 8 页为空白，译本不另造空白页。

\SGAfrontchapter{引言}
\SourcePage{9}
这里呈献给读者的是高等科学研究所于 1962 年举行的第二期\Term{algebra}{代数}几何讨论班讲义（油印本）的修订增补版，以照相胶印方式重新出版。

关于这些讨论班所追求的目标及其与《\Term{algebra}{代数}几何基础》（EGA）的关系，请读者参阅第一期讨论班（下文简称 SGA 1）的引言。

第 I 至 XI 讲的文字由一组听众根据我的口头报告和手写笔记陆续整理而成，这些听众包括 I. Giorgiutti、J. Giraud、M. Jaffe 女士（后来改姓为 M. Hakim）和 A. Laudal。这些笔记起初被视为临时性的材料，只在极小范围内传阅，等待日后收入 EGA（如今至少可以说，这一收入计划已经很成问题，SGA 的其他部分也是如此）。正如初版的说明所说，笔记的这种“内部”性质，应当可以使其中某些“行文上的不足”得到谅解；这些不足在本期 SGA 2 中无疑比其他各期更为明显。在本次重版中，我尽可能通过对原文作较为细致的修订来弥补这些不足。特别是，我统一了各讲陈述的编号方式，通篇采用同一种十进编号法；本期 SGA 2 的初版多数讲以及 SGA 的其他各部分，原已采用这一编号法。因此，我尤其不得不全面改动第 III 至 VIII 讲中各陈述的编号\footnote[1]{编者注：我们尽可能保留了原有编号；对于偶尔出现的、有歧义的重复编号，则加上副词 \textit{bis}。}，因而也改动了对这些讲的引用。\SourcePage{10}{\renewcommand{\thefootnote}{\fnsymbol{footnote}}\footnote[1]{不言而喻，在 \textit{Series in Pure Mathematics} 丛书所出版的 SGA 各部分中，一切对 SGA 2 的引用都将指向本卷，而不是 SGA 2 的初版！}}我也试图清除原文中主要的打字错误和语法错误（这些错误数量很多，且妨碍阅读）。此外，M. Hakim 女士欣然承担了重写第 IV 讲的工作，使其行文不像原讲义那样简略。如同 SGA 的其他各次重版，我还增添了若干页下注，或提供补充参考文献，或说明某个问题自原文写成以来所取得进展的状况。最后，本期讨论班还增添了一讲，即 Michèle Raynaud 女士于 1967 年写成的第 XIV 讲；它重新讨论并补充了“第 XIII 讲评注”（XIII 6，写于 1963 年 3 月）中的一些建议。这一讲从\Term{etale-cohomology}{平展上同调}的观点重新讨论 Lefschetz 型定理，使用了 SGA 4 和 SGA 5 所阐述的\Term{etale-cohomology}{平展上同调}结果（这两部书将由同一套 \textit{Series in Pure Mathematics} 丛书出版）\footnote[2]{编者注：实际上，这些讨论班讲义由 Springer 出版，卷号为 269、279、305 和 589，但遗憾的是，均已售罄。}；因此，它的性质不如本卷其他各讲那样“初等”，后者所使用的材料几乎没有超出 EGA 第 I 至 III 章的内容。下面概述本卷内容。第 I 讲给出了“\Term{cohomology-with-support}{支集在 $Y$ 中的上同调}”$\mathrm{H}^{*}_{Y}(X,F)$ 的一系列基本性质，其中 $Y$ 是空间 $X$ 的一个闭集。这种\Term{cohomology}{上同调}可解释为 $X$ 模去\Term{open-subset}{开集} $X-Y$ 的\Term{cohomology}{上同调}；它是一个很有用的“\Term{local-to-global-spectral-sequence}{从局部到整体的谱序列}”I 2.6 的收敛目标，该\Term{spectral-sequence}{谱序列}涉及“\Term{support}{支集}在 $Y$ 中”的\Term{cohomology-sheaf}{上同调层} $\underline{\mathrm{H}}^{*}_{Y}(F)$\footnote[3]{编者注：我们保留了用下划线表示\Term{functor}{函子}的\Term{sheafification}{层化}版本的记号，因为 $\Gamma$ 对应的书写体记号并不明确。}。在许多问题中，这套形式体系可以起到“\Term{localization}{局部化}”的\Term{action}{作用}，类似于\Term{differential}{微分}几何中考虑 $Y$ 的“管状”邻域所起的\Term{action}{作用}。第 II 讲在\Term{prescheme}{预概形}上的\Term{quasi-coherent-sheaf}{拟凝聚层}这一情形下研究上述概念；第 III 讲给出它们与经典的\Term{depth}{深度}概念之间的关系（III 3.3）。

第 IV、V 讲给出\Term{local-duality}{局部对偶}的概念，可以将其与 Serre 的\Term{projective-duality}{射影对偶}定理（XII 1.1）相比较。应当指出，这两类\Term{duality}{对偶}定理在 Hartshorne 的讨论班讲义中都得到了实质性的推广（该讲义引于第 IV 讲末的页下注）\footnote[4]{编者注：Hartshorne 的书中有符号错误，尤其没有真正证明迹与\Term{base-change}{基变换}的相容性。Conrad 完全重新处理了这项工作，证明了这种至关重要而且极不\Term{trivial-bundle}{平凡}的相容性（\textit{Grothendieck duality and base change}，Lecture Notes in Mathematics，第 1750 卷，Springer-Verlag，Berlin，2000）。遗憾的是，仍有错误存在（参见 Conrad B. 的两篇预印本：“Clarifications and corrections to `Grothendieck duality and base change'”和“An addendum to Chapter 5 of `Grothendieck duality and base change'”）。关于更具体的方面，尤其是\Term{residue}{留数}概念，参见 Lipman 的工作，特别是 Lipman J.，\textit{Dualizing sheaves, differentials and residues on algebraic varieties}，Astérisque，第 117 卷，Société mathématique de France，1984。一种基于 Brown \Term{representable}{可表}性定理的\Term{duality}{对偶}定理的\Term{category}{范畴}论证明由 Neeman 得到：Neeman A.，“The Grothendieck duality theorem via Bousfield's techniques and Brown representability”，\textit{J. Amer. Math. Soc.} 9（1996），第 1 期，第 205--236 页。}。

\SourcePage{11}
第 VI、VII 讲给出一些容易的技术性概念，第 VIII 讲利用它们证明\Term{finiteness-theorem}{有限性定理}（VIII 2.3）。该定理给出如下必要充分条件：对于\Term{noetherian-scheme}{诺特概形} $X$ 上的\Term{coherent-sheaf}{凝聚层} $F$，其\Term{local-cohomology-sheaf}{局部上同调层} $\underline{\mathrm{H}}^{i}_{Y}(F)$ 在 $i\leq n$ 时为\Term{coherent-sheaf}{凝聚层}；\Term{equivalence}{等价}地，层 $\mathrm{R}^{i}f_{*}(F|_{X-Y})$ 在 $i\leq n-1$ 时为\Term{coherent-sheaf}{凝聚层}，这里 $f\colon X-Y\longrightarrow X$ 为包含映射。该定理是本期讨论班技术方面的核心结果之一。我们在第 IX 讲中说明，如何仿照并推广 EGA III 第 4、5 节中对\Term{proper-morphism}{固有态射}的\Term{finiteness-theorem}{有限性定理}的运用，利用这类定理在\Term{formal-geometry}{形式几何}中建立“\Term{comparison-theorem}{比较定理}”和“\Term{existence-theorem}{存在定理}”。

第 X、XI 讲分别讨论\Term{fundamental-group}{基本群}和 \Term{picard-group}{Picard 群}的 Lefschetz 型定理，并应用上述结果。这些定理在一定条件下，比较分别附属于\Term{scheme}{概形} $X$ 及其一个\Term{subscheme}{子概形} $Y$（起\Term{hyperplane-section}{超平面截面}的\Term{action}{作用}）的\Term{invariants}{不变量}（$\pi_{1}$ 或 $\operatorname{Pic}$），尤其给出它们\Term{isomorphism}{同构}的条件。大体说来，所作假设用以从 $Y$ 过渡到 $X$ 沿 $Y$ 的\Term{formal-completion}{形式完备化}，然后应用第 IX 讲的结果，从那里过渡到 $X$ 中 $Y$ 的一个\Term{open-neighborhood}{开邻域} $U$。为了从 $U$ 过渡到 $X$，还需要掌握 $X$ 在 $Z=X-U$ 的各点处的\Term{local-ring}{局部环}的进一步信息（“\Term{purity}{纯性}”或“\Term{parafactoriality}{副析因性}”一类的信息）；在所考虑的情形中，$Z$ 是一个\Term{finite}{有限}\Term{discrete}{离散}集。这就说明了第 X、XI、XII 讲的证明中\Term{local-scheme}{局部}结果与整体结果之间的相互\Term{action}{作用}，尤其在某些归纳论证中的相互\Term{action}{作用}。第 X、XI 讲所得主要结果是具有\Term{local-scheme}{局部}性质的定理 X 3.4（\Term{purity-theorem}{纯性定理}）和 XI 3.14（\Term{parafactoriality-theorem}{副析因性定理}）。应当注意，这些定理是由本质上具有整体性质的\Term{cohomology}{上同调}方法证明的。第 XII 讲运用上述\Term{local-scheme}{局部}结果，得到这些结果对于域上的\Term{projective-scheme}{射影概形}，或更一般地，对于一个或多或少任意的\Term{base-scheme}{基概形}上的\Term{projective-scheme}{射影概形}的整体变体；作为典型的陈述，特别指出 XII 3.5 和 XII 3.7。

\SourcePage{12}
第 XIII 讲回顾由本期讨论班的结果和方法引出的众多问题与猜想中的一部分。其中最有趣的也许是有关\Term{complex-analytic-space}{复解析空间}的\Term{cohomology}{上同调}与\Term{homotopy}{同伦}的 Lefschetz 型定理，参见第 XIII 讲第 26 页及以下各页\footnote[5]{编者注：第 XIII、XIV 讲提出的猜想如今基本上都已得到证明；参考文献与评注见这两讲的页下注。}。在\Term{scheme}{概形}的\Term{etale-cohomology}{平展上同调}框架中，相应的猜想在第 XIV 讲借助一种\Term{duality}{对偶}技术得到证明，这种技术应当同样适用于复解析情形（参见第 XIII 讲评注第 25 页和 XIV 6.4）。但是，在\Term{analytic-space}{解析空间}的情形下，相应的\Term{homotopy}{同伦}论陈述（尤其是涉及\Term{fundamental-group}{基本群}的陈述）似乎需要全新的技术（参见 XIV 6.4）。

我很高兴在此感谢所有以不同方式帮助本卷出版的人，包括本引言中已经提到的合作者。尤其要感谢 Chardon 女士，她欣然承担了为照相胶印准备最终书稿这一繁琐而少人乐做的工作。

\begin{flushright}
Bures-sur-Yvette，1968 年 4 月\\
A. Grothendieck
\end{flushright}
